%表紙
\newpage
\chapter{運営・組織}
\begin{enumerate}[label=第\arabic*条　]
  \item この章は、憲章の趣旨に基づき、学園の運営に必要な役職及びその選任等に関し必要な事項を定めるものとする。
  \item 学園に、学園長及び教務主事を置く。
    \begin{enumerate}[label=\arabic*　,start=2,leftmargin=*]
      \item 前項に定めるもののほか、学園に講師その他の役職を置くことができる。
    \end{enumerate}
  \item 教務主事となろうとする学生は、学園長に申し出るものとする。
    \begin{enumerate}[label=\arabic*　,start=2,leftmargin=*]
      \item 学園長は、前項の申出を承認したときは、当該学生を教務主事に委嘱する。
      \item 教務主事の人数について、定数は定めない。
    \end{enumerate}
  \item 学園長は、教務主事のうちの一人をもって充てる。
    \begin{enumerate}[label=\arabic*　,start=2,leftmargin=*]
      \item 学園長が欠けたときは、教務主事の協議により、その後任を定める。
    \end{enumerate}
  \item 教務主事の協議が調わないとき、又は緊急を要するときは、学園長が決定することができる。
    \begin{enumerate}[label=\arabic*　,start=2,leftmargin=*]
      \item 前項の規定は、次の各号に掲げる事項については、適用しない。
        \begin{enumerate}[label=(\arabic*)　]
          \item 第６章第２条及び第７章第６条の規定による認可
          \item 学園長の解任に関する事項
        \end{enumerate}
      \item 第１項の規定による決定は、教務主事の間の意思決定を円滑に行うための手続であって、学園長が他の教務主事に優越することを意味するものではない。
    \end{enumerate}
  \item 役職者の活動は、無報酬とする。役職者は、その職務の遂行に対し、報酬その他の対価を受けないものとする。
    \begin{enumerate}[label=\arabic*　,start=2,leftmargin=*]
      \item 前項の規定は、職務の遂行に対する対価としての性質を有しない、儀礼の範囲内の贈与を受けることを妨げない。
    \end{enumerate}
  \item 役職者は、学生としての資格を失ったときは、当然に役職者としての地位を失う。
  \item 前条に定めるもののほか、役職者としての地位の喪失に関する事項は、第８章に定めるところによる。
  \item この章の実施に関し必要な事項は、教務主事が別に定める。
\end{enumerate}
